%% kapa: exam=1 round=37 year=2026 period=2 form=A subject=appraisal_law subject_ko=감정평가관계법규 seq=35 number=35
%% source: Q-Net 감정평가사 기출문제 (cst00309) artlSeq=5258924 file=raw/1008/5258924/2026년 제37회 감정평가사 1차 2교시.pdf
%% conversion: kordoc 4.24.51 -> tools/gen1.py (auto); review=proofread-sonnet; flags=none
\begin{question}{35}
부동산등기법령상 등기절차에 관한 설명으로 옳지 않은 것은?
\begin{choices}
  \item 토지의 합병이 있는 경우에는 그 토지 소유권의 등기명의인은 그 사실이 있는 때부터 1개월 이내에 그 등기를 신청하여야 한다.
  \item 신탁재산에 속하는 부동산의 신탁등기는 수탁자가 단독으로 신청한다.
  \item 등기관이 등기를 마친 후 그 등기에 착오가 있음을 등기권리자, 등기의무자 또는 등기명의인에게 알려야 하는 경우 이들이 각 2인 이상이면 그 모두에게 통지하여야 한다.
  \item 등기관이 등기의 빠진 부분이 등기관의 잘못으로 인한 것임을 발견한 경우 지체 없이 그 등기를 직권으로 경정하여야 하며, 등기상 이해관계 있는 제3자가 있다면 제3자의 승낙이 있어야 한다.
  \item 대표자나 관리인이 있는 법인 아닌 사단이나 재단에 속하는 부동산의 등기에 관하여는 그 사단이나 재단을 등기권리자 또는 등기의무자로 한다.
\end{choices}
\end{question}
\answer{3}
